\subsection{Examples of measures}

Example 1. --- Measures on a locally compact space.

The following proposition shows that the theory of this chapter contains that of Ch. IV. In the statement, the word `measure' and the notation \(\mathcal{M}(\mathrm{T};\mathbf{C})\) are taken in the sense of the earlier chapters.

PROPOSITION 3. --- Let T be a locally compact space, and let \(\mu\) be a measure on T. Denote by W( \(\mu\) ) the mapping that associates to each compact subset K of T the induced measure \(\mu_{K}\) . Then W( \(\mu\) ) is a premeasure on T, one has W(\textbar{} \(\mu\) \textbar) = \textbar W( \(\mu\) )\textbar, and the linear mapping W: \(\mu \mapsto\) W( \(\mu\) ) is a bijection of the space \(\mathcal{M}(T;C)\) onto the space \(\mathcal{P}(T;C)\) of premeasures on T. Moreover, if \(\mu\) is positive then \(\mu^{\bullet} = (W(\mu))^{\bullet}\) .

It is clear that \(\mathrm{W}(\mu)\) is a premeasure (Ch. V, §7, No.~2, Prop. 4) and that the mapping W is linear. The relation \(\mathrm{W}(\mu)=0\) means that \(\mu\) induces the measure 0 on every compact set in T; then \(\mu(f)=0\) for \(f\in\mathcal{K}(\mathrm{T};\mathbf{C})\) , thus \(\mu=0\) , which proves that W is injective. It remains to prove that W is surjective. Since every premeasure is a linear combination of positive premeasures, it will suffice to construct, for every positive premeasure w, a positive measure \(\mu\) such that \(w=\mathrm{W}(\mu)\) . Let \(f\in\mathcal{K}(\mathrm{T})\) be a given function, and let L be a compact set containing the support of f; the number \(w_{\mathrm{L}}(f_{\mathrm{L}})\) is independent of the choice of L, by the definition of induced measures, so that one can set \(\mu(f)=w_{\mathrm{L}}(f_{\mathrm{L}})\) ; then \(\mu\) is a positive linear form on \(\mathcal{K}(\mathrm{T})\) , that is, a positive measure. Let us verify that \(w = \mathrm{W}(\mu)\) ; first, the relation \(\mu^{\bullet}(f) = w_{\mathrm{L}}^{\bullet}(f_{\mathrm{L}})\) extends to the case that \(f\) is a finite upper semi-continuous function that is positive and is zero outside \(\mathbf{L}\) . For, let \(\mathbf{M}\) be a compact neighborhood of \(\mathbf{L}\) , and \(\mathcal{H}\) the (decreasing directed) set of continuous functions on \(\mathbf{T}\) , with support contained in \(\mathbf{M}\) , that are \(\geqslant f\) . Then (Ch. IV, §4, No.~4, Cor. 2 of Prop. 5)

\[
\mu^ {\bullet} (f) = \inf _ {h \in \mathscr {H}} \mu (h) = \inf _ {h \in \mathscr {H}} w _ {\mathrm{M}} (h _ {\mathrm{M}}) = w _ {\mathrm{M}} ^ {\bullet} (f _ {\mathrm{M}}),
\]

and on the other hand \(w_{\mathrm{M}}^{\bullet}(f_{\mathrm{M}}) = w_{\mathrm{L}}^{\bullet}(f_{\mathrm{L}})\) since \(f_{M}\) is zero on M-L (Ch. V, §7, No.~1, Prop. 1). In particular, if f is taken to be the extension by 0 of an element of \(\mathcal{K}_{+}(\mathrm{L})\) , this formula shows that \(\mu_{L} = w_{L}\) by the definition of induced measures, thus indeed \(\mathrm{W}(\mu) = w\) .

If \(\mu\) is positive, then

\[
\mu^ {\bullet} (f) = \sup _ {K} \mu^ {\bullet} (f \varphi_ {K}) = \sup _ {K} \mu_ {K} ^ {\bullet} (f _ {K}) = (W (\mu)) ^ {\bullet} (f)
\]

for every \(f \in \mathcal{F}_{+}(\mathrm{T})\) (Ch. V, §1, Def. 1 and §7, Prop. 1). The relation \(|\mathrm{W}(\mu)| = \mathrm{W}(|\mu|)\) is obvious (Ch. IV, §5, No.~7, Lemma 3).

Q.E.D.

When T is locally compact, we shall from now on identify the spaces \(\mathcal{M}(\mathrm{T};\mathbf{C})\) and \(\mathcal{P}(\mathrm{T};\mathbf{C})\) by means of the bijection W.

Example 2. --- Measures with compact support on a topological space.

Lemma 2. --- Let T be a topological space, L a compact subset of T, and \(\lambda\) a positive measure on L. There exists a unique positive measure \(\mu\) on T such that, for every function \(f \in \mathcal{F}_{+}(\mathrm{T})\) ,

\[
\mu^ {\bullet} (f) = \lambda^ {\bullet} (f _ {L}).\tag{2}
\]

Let us set \(p(f) = \lambda^{\bullet}(f_{\mathrm{L}})\) for every \(f \in \mathcal{F}_{+}(\mathrm{T})\) , and let us show that the conditions 1) and 2) of Prop. 2 b) are satisfied. The second is obviously satisfied: indeed, \(p(f) = p(f\varphi_{\mathrm{K}})\) if \(\mathrm{K}\) contains \(\mathrm{L}\) . If \(\mathrm{K} \subset \mathrm{T}\) is compact, and if \(h \in \mathcal{F}_{+}(\mathrm{K})\) , then

\[
p _ {\mathrm{K}} (h) = p \left(h ^ {0}\right) = \lambda^ {\bullet} \left(h ^ {0} \mid \mathrm{L}\right).
\]

But \(h^0 | \mathbf{L}\) is the extension by 0 of \(h_{\mathbf{K} \cap \mathbf{L}}\) to \(\mathbf{L}\) : the last expression is therefore equal to \((\mu_{\mathbf{K}})^{\bullet}(h)\) , where \(\mu_{\mathbf{K}}\) is the image of \(\lambda | \mathbf{K} \cap \mathbf{L}\) under the injection of \(\mathbf{K} \cap \mathbf{L}\) into \(\mathbf{K}\) (Ch. V, §6, No.~2, Prop. 2 and §7, No.~1, Prop. 1), and \(p_{\mathbf{K}} = (\mu_{\mathbf{K}})^{\bullet}\) . Condition 1) of Prop. 2 b) is therefore also satisfied, and the existence of \(\mu\) follows at once.

Q.E.D.

We shall say that \(\mu\) is the measure on T defined by \(\lambda\) . In particular, for every point \(x\) of T one can define the measure \(\varepsilon_x\) ; it is characterized by \((\varepsilon_x)^{\bullet}(f) = f(x)\) for \(f \in \mathcal{F}_{+}(\mathrm{T})\) .

Remarks. --- 1) When T is locally compact, \(\mu\) is the image of \(\lambda\) under the injection of L into T. We shall see in §2, No.~3, Example, when image measures will have been treated, that this interpretation remains valid for arbitrary spaces.

\begin{enumerate}
\def\labelenumi{\arabic{enumi})}
\setcounter{enumi}{1}
\tightlist
\item
  We shall also see that the measures defined in Example 2 are positive measures on T with compact support (No.~6, Remark 2).
\end{enumerate}

We shall henceforth consider only positive measures, absent express mention to the contrary. For the rest of this section, T will denote a topological space and \(\mu\) a positive measure on T.

Numerous results in the following subsections may be extended to positive premeasures. This extension is left to the reader.

